% from ega2-1-fr.tex
\oldpage[II]{139}

il existe alors un anneau de valuation ayant $K$ pour corps des fractions et
qui domine $A$ (\cite[p.~1-07, lemme~2]{I-1}).

D'autre part, soient $B$ un anneau de valuation, $k$ son corps des restes,
$K$ son corps des fractions, $L$ une extension de $k$. Alors il existe un
anneau de valuation \emph{complet} $C$ qui domine $B$ et dont le corps des
restes est $L$. En effet, $L$ est extension algébrique d'une extension
transcendante pure $L'=k(T_\mu)_{\mu\in M}$~; on sait qu'on peut prolonger la
valuation de $K$ correspondant à $B$ à une valuation de
$K'=K(T_\mu)_{\mu\in M}$ de sorte que $L'$ soit le corps des restes de cette
valuation (\cite[p.~98]{II-24})~; remplaçant $B$ par le complété de l'anneau
de cette valuation prolongée, on voit qu'on peut se limiter au cas où $B$ est
complet et $L$ est une clôture algébrique de $k$. Si $\overline{K}$ est une
clôture algébrique de $K$, on peut alors prolonger à $\overline{K}$ la
valuation qui définit $B$, et le corps des restes correspondant est une
clôture algébrique de $k$, comme on le voit en relevant dans $\overline{K}$
les coefficients d'un polynôme unitaire de $k[T]$. On est donc finalement
ramené au cas où $L=k$ et il suffit alors de prendre pour $C$ le complété de
$B$ pour répondre à la question.
\end{env}

\begin{env}[7.1.3]
\label{II.7.1.3-fr}
Soient $K$ un corps, $A$ un sous-anneau de $K$~; la fermeture intégrale $A'$
de $A$ dans $K$ est l'intersection des anneaux de valuation contenant $A$ et
ayant $K$ pour corps des fractions
(\cite[p.~51, th.~2]{I-11}). La proposition
(\hyperref[II.7.1.2-fr]{7.1.2}) s'exprime sous la forme géométrique
équivalente~:
\end{env}

\begin{proposition}[7.1.4]
\label{II.7.1.4-fr}
Soient $Y$ un préschéma, $p:X\to Y$ un morphisme, $x$ un point de $X$,
$y=p(x)$, $y'\neq y$ une spécialisation
(\hyperref[0.2.1.2-fr]{0, 2.1.2}) de $y$. Il existe alors un schéma local
$Y'$, spectre d'un anneau de valuation, et un morphisme séparé $f:Y'\to Y$
tel que, si $a$ désigne l'unique point fermé de $Y'$ et $b$ le point
générique de $Y'$, on ait $f(a)=y'$ et $f(b)=y$. On peut en outre supposer
vérifiée l'une des deux propriétés additionnelles suivantes~:
\begin{enumerate}
\item[\rm{(i)}] $Y'$ est le spectre d'un anneau de valuation \emph{complet}
dont le corps des restes est algébriquement clos, et il existe un
$k(y)$-homomorphisme $k(x)\to k(b)$.
\item[\rm{(ii)}] Il existe un $k(y)$-isomorphisme
$k(x)\xrightarrow{\sim}k(b)$.
\end{enumerate}
\end{proposition}

Soit $Y_1$ le sous-préschéma fermé réduit de $Y$ ayant
$\overline{\{y\}}$ comme espace sous-jacent
(\hyperref[I.5.2.1-fr]{I, 5.2.1}), et soit $X_1$ le sous-préschéma fermé
image réciproque $p^{-1}(Y_1)$~; comme $y'\in\overline{\{y\}}$ par hypothèse
et que $k(x)$ est le même dans $X$ et dans $X_1$, on peut supposer $Y$
\emph{intègre}, de point générique $y$~; $\mathcal{O}_{y'}$ est alors un
anneau local intègre qui n'est pas un corps, et dont le corps des fractions
est $\mathcal{O}_y=k(y)$, et $k(x)$ est une extension de $k(y)$. Pour
réaliser les conditions $f(a)=y'$ et $f(b)=y$ ainsi que la condition
additionnelle (i) (resp. (ii)), on prendra
$Y'=\operatorname{Spec}(A')$, où $A'$ est un anneau de valuation dominant
$\mathcal{O}_{y'}$ et qui est complet et a un corps des restes
algébriquement clos extension de $k(x)$ (resp. un anneau de valuation $A'$
dominant $\mathcal{O}_{y'}$ et dont $k(x)$ est le corps des fractions)~;
l'existence de $A'$ est garantie par
(\hyperref[II.7.1.2-fr]{7.1.2}).

\begin{env}[7.1.5]
\label{II.7.1.5-fr}
Rappelons qu'un anneau local $A$ est dit \emph{de dimension $1$} s'il existe
un idéal premier distinct de l'idéal maximal $\mathfrak{m}$, et si tout idéal
premier de $A$ distinct de $\mathfrak{m}$ est un idéal premier
\emph{minimal}~; lorsque $A$ est \emph{intègre}, il revient au même de dire
que $\mathfrak{m}$ et $(0)$ sont les seuls idéaux premiers et
$\mathfrak{m}\neq(0)$~; autrement dit, $Y=\operatorname{Spec}(A)$ est
réduit à deux

